\section{Lifting Idempotents}

\begin{definition}
    Let $R$ be a ring and let $I$ be a two-sided ideal of $R$.

    If $f\in R$ is an idempotent and $ f \notin I $, then $e=f+I$ is an idempotent of $R/I$. We say that an idempotent $f\in R$ \emph{lifts}
    $e$.
\end{definition}

\begin{remark}
    For a general ideal $I$, an idempotent of $R/I$ need not lift to an idempotent of $R$.
    The next theorem shows that lifting is always possible when $I$ is nilpotent.
\end{remark}

\begin{theorem}\label{thm:lift-idempotent}
    Let $R$ be a ring, let $I$ be a nilpotent two-sided ideal of $R$, and let
    \[
        e\in R/I
    \]
    be an idempotent. Then there exists an idempotent $f\in R$ such that
    \[
        f+I=e.
    \]
    Moreover, if $e$ is primitive, then every lift $f$ of $e$ is primitive.
\end{theorem}

\begin{proof}
    Since $I$ is nilpotent, there exists $r\ge 1$ such that
    \[
        I^r=0.
    \]
    We shall construct idempotents
    \[
        e_i\in R/I^i \qquad (1\le i\le r)
    \]
    such that
    \[
        e_1=e
    \]
    and for each $i\ge 2$, the image of $e_i$ in $R/I^{i-1}$ is $e_{i-1}$.

    Assume inductively that $e_{i-1}\in R/I^{i-1}$ is an idempotent, with $2\le i\le r$.
    Choose any element
    \[
        a\in R/I^i
    \]
    whose image in $R/I^{i-1}$ is $e_{i-1}$.
    Since $e_{i-1}$ is idempotent, we have
    \[
        a^2-a\in I^{i-1}/I^i.
    \]
    Now
    \[
        \bigl(I^{i-1}/I^i\bigr)^2=0,
    \]
    because
    \[
        I^{2i-2}\subseteq I^i
        \qquad (i\ge 2).
    \]
    Hence
    \[
        (a^2-a)^2=0
        \qquad\text{in } R/I^i.
    \]

    Define
    \[
        e_i:=3a^2-2a^3\in R/I^i.
    \]
    First, $e_i$ maps to $e_{i-1}$ in $R/I^{i-1}$, because
    \[
        e_i-a
        =
        3a^2-2a^3-a
        =
        -(2a-1)(a^2-a)\in I^{i-1}/I^i.
    \]
    Next, $e_i$ is idempotent. Indeed, the polynomial identity
    \[
        (3t^2-2t^3)^2-(3t^2-2t^3)
        =
        -(3-2t)(1+2t)(t^2-t)^2
    \]
    gives
    \[
        e_i^2-e_i
        =
        -(3-2a)(1+2a)(a^2-a)^2=0
    \]
    in $R/I^i$.

    Thus, by induction, we obtain an idempotent
    \[
        e_r\in R/I^r=R.
    \]
    Let
    \[
        f:=e_r.
    \]
    Then $f$ is an idempotent of $R$ whose image in $R/I$ is $e$.

    \medskip

    Now suppose that $e$ is primitive and that $f$ is a lift of $e$.
    Assume that
    \[
        f=f_1+f_2
    \]
    is a decomposition of $f$ into orthogonal idempotents:
    \[
        f_1^2=f_1,\qquad f_2^2=f_2,\qquad f_1f_2=f_2f_1=0.
    \]
    Reducing modulo $I$, we get
    \[
        e=(f_1+I)+(f_2+I),
    \]
    and this is again a decomposition into orthogonal idempotents in $R/I$.
    Since $e$ is primitive, one of $f_1+I$ or $f_2+I$ must be zero.
    Without loss of generality, assume
    \[
        f_1+I=0.
    \]
    Then $f_1\in I$.

    But $f_1$ is idempotent, so for every $m\ge 1$,
    \[
        f_1=f_1^m\in I^m.
    \]
    Taking $m=r$, we get
    \[
        f_1\in I^r=0,
    \]
    hence $f_1=0$.
    Therefore $f$ cannot be written as a sum of two nonzero orthogonal idempotents, so $f$ is primitive.
\end{proof}

